\documentclass[conference,a4paper]{IEEEtran}
\IEEEoverridecommandlockouts
\usepackage{cite}
\usepackage{url}
\usepackage{amsmath,amssymb,amsfonts}
\usepackage{algorithmic}
\usepackage{graphicx}
\usepackage{textcomp}
\usepackage{xcolor}
\usepackage{booktabs}
\def\BibTeX{{\rm B\kern-.05em{\sc i\kern-.025em b}\kern-.08em
    T\kern-.1667em\lower.7ex\hbox{E}\kern-.125emX}}
\begin{document}

\title{Hand Gesture Mouse Control for Lean-Back Computing}

\author{\IEEEauthorblockN{Leo}
\textit{Victoria University of Wellington}\\
AIML339 Final Project Report}

\maketitle

\begin{abstract}
We ask whether a decision-and-verification layer around an existing palm
detector can drive a screen cursor by hand at 2 to 4 m in dim artificial
light, where MediaPipe Hands degrades. We adopt a third-party PyTorch port
of BlazePalm, train a 60,929-parameter crop verifier to replace a hand-tuned
size gate, add a classical motion gate, and fine-tune 2.49\% of the
detector's parameters on about 200 dark frames. Scoring all seven systems in
one session against identical ground truth, the shipped system reaches F1
50.3\% against MediaPipe's 81.6\%, with no false locks on empty-room footage
for any configuration. Paired tests support that gap: MediaPipe decides
correctly on 553 frames where the shipped system does not against 38 the
other way (Holm-adjusted $p<10^{-100}$), with an F1 difference of $+0.31$
whose 95\% block-bootstrap interval is $+0.23$ to $+0.40$. At close range F1
is 85 to 94\% against 100\%. A dark-specialised variant raises
VeryLowLight F1 from 10.4\% to 31.5\%, but that clip was partly used for
fine-tuning and the variant collapses on close static hands (93.8\% to
9.4\%). A staged analysis shows the palm detector loses 64.3\% of
target-condition hand frames, which is 86\% of all frames lost once the
verifier is fixed. Negative results: two classical contrast baselines in
front of unmodified MediaPipe (F1 80.4 with CLAHE, 79.8 with gamma, against
81.6 raw), a logistic baseline worse than no filtering, and a motion-blob
path with a 4.8-point ceiling. Verifier variants were retrained with 3 to 10
seeds each; held-out AUC is stable to within 0.008, while the operating
threshold chosen on 151 validation frames is not.
\end{abstract}

\begin{IEEEkeywords}
hand tracking, gesture interaction, MediaPipe, BlazePalm, learned verification, low-light detection
\end{IEEEkeywords}

\section{Introduction}

The target is a gesture mouse for lean-back computing: a webcam tracks a
hand 2 to 4 m from the screen, in dim artificial light, and the hand drives
the cursor. Three measured properties define the problem. First, range: the
palm box is about 0.085 of the frame width (136 px at 1920 wide, 116 px at
1280). Second, lighting: our target-condition recording, VeryLowLight, has
median luma 43.5 against 29 to 146 for the other clips, contrast (standard
deviation) near 17 against 36 to 53, and a colour cast R/B of 1.79 against
1.05 to 1.21. MediaPipe detects a hand in only 60.8\% of its frames in
tracking mode and 44.7\% in static mode on this clip. Third, latency
tolerance is asymmetric: slow, deliberate movement is fine, but a cursor
that jumps to the wrong place is not.

\textbf{Research question.} Does a decision-and-verification layer around a
MediaPipe-class palm detector beat MediaPipe at range and in low light, and
where does the remaining error come from?

\textbf{Change from the proposal.} The proposal was to modify MediaPipe's
tracking architecture and add an open/closed hand-state output. We
abandoned that for a PyTorch port of BlazePalm, because none of the
components below can be inserted into a compiled MediaPipe graph. The
open/closed output was dropped when the landmark model's presence score
fell to a median of 0.001 on VeryLowLight at locations MediaPipe itself
labels as a hand (1.000 in daylight). The final system reports hand
location and lock state only.

\textbf{Contribution.} This is not a new hand-tracking architecture. It is
(i) a staged attribution of where a detector-based pipeline fails at range
and in low light, (ii) a set of small learned and classical components,
each measured on its own, and (iii) several documented failures.
Section~\ref{sec:method} separates reused, modified, extended and new work.

\section{Background and Related Work}

\textbf{MediaPipe Hands and BlazePalm.} MediaPipe Hands \cite{mediapipe}
uses a palm detector to propose a box, then a landmark model regresses 21
keypoints in a rotation-normalised crop. The detector is a single-shot
regressor in the BlazeFace family \cite{blazeface}, with 2944 unit-size
anchors and six 1$\times$1 heads, released in the MediaPipe graph
\cite{mediapipe_github}. We use the PyTorch reimplementation of Satija
\cite{blazepalm_port}, which loads Google's released weights and is not
retrained. Because its anchors are all unit-size, MediaPipe's scale-aware
anchor matching does not carry over, which matters for fine-tuning
(Section~\ref{sec:method}).

\textbf{Data.} HaGRID \cite{hagrid} and HaGRIDv2 \cite{hagrid2} contain
hands recorded at 0.5 to 4 m under varied light, which matches our range
but has no true near-dark footage. We pulled HaGRIDv2 from two Hugging Face
mirrors (Section~\ref{sec:ethics}).

\textbf{Baseline choice.} MediaPipe is both the method we try to improve on
and, through its landmark presence head, our label arbiter. Its detection
output also defines the ground truth, so every result here measures
agreement with MediaPipe, not with a human annotator. We also test the
cheap rival that a low-light report is expected to rule out: classical
contrast enhancement in front of unmodified MediaPipe, specifically CLAHE
\cite{clahe} on the L channel of LAB (clip limit 2.0, $8\times8$ tiles) and
luminance gamma correction at $\gamma=0.5$. Both are applied with textbook
parameters fixed in advance, not tuned on any part of our data. Neither
helps (Section~\ref{sec:results}), which is a negative result we report
rather than omit: it means the target-condition loss is not a contrast
problem that a preprocessing step can solve.

\section{Methodology and Implementation}
\label{sec:method}

\subsection{What was reused, modified, extended and built}
Reused: the BlazePalm architecture and its released weights (1,763,358
parameters), MediaPipe as the baseline, the ground truth and the label
arbiter, and the PyTorch and OpenCV libraries. Modified: the six
$1\times1$ prediction heads (43,966 parameters, 2.49\%), fine-tuned for
the dark variant only, with the original weights kept byte-identical.
Extended: the acquisition motion gate together with the relaxed size gate
it makes safe, the motion-blob proposals, and the two consensus rules.
Built here: the CNN verifier and its data pipeline, the logistic baseline,
the state machine with its plausibility gate, smoothing and cursor control,
and the evaluation harness with its arbiter-based label cleaning.

\begin{figure}[t]
\centerline{\includegraphics[width=0.49\textwidth]{fig1_pipeline.png}}
\caption{Pipeline. The verifier confirms every candidate from the detector
or the blob path before it can move the cursor. The motion gate acts during
acquisition only.}
\label{fig:pipeline}
\end{figure}

\subsection{Palm detector and fine-tune}
The detector takes a 256$\times$256 letterboxed frame scaled to $[-1,1]$
and decodes boxes from 2944 anchors. The port discards scores below 0.7
before weighted NMS, so the application's own 0.5 floor never binds. For
the dark variant only, we freeze everything except the six 1$\times$1 heads
and fine-tune on about 200 VeryLowLight frames (Adam, lr $3\times10^{-4}$,
batch 4, 20 epochs, seed 0). The loss is
$\mathcal{L}_{\text{cls}} + 0.5\,\mathcal{L}_{\text{reg}}$, with positive
weight 5.0 and a masked smooth-L1 over the box and two of the seven
keypoints. Labelling every anchor inside the ground-truth box (about 96 of
2944) positive destroyed the score head, firing on 337 anchors instead of
13, so the shipped version matches only the 6 anchors nearest the box centre
and a guard refuses to save a model firing on more than 100 anchors.

\subsection{Crop verifier}
A small CNN (four stride-2 Conv--BatchNorm--ReLU blocks,
3$\to$16$\to$32$\to$64$\to$64 channels, global average pooling, one linear
output; 60,929 parameters) takes the 64$\times$64 RGB rotation-normalised
crop the pipeline already computes and outputs $P(\text{hand})$. It
replaces a box-size gate that was the main source of false locks at range.
Training uses BCE, Adam, lr $2\times10^{-3}$, batch 64, 8 epochs, on a
class-balanced subsample. Four variants were shipped, trained on 11k, 130k
(arbiter-filtered), 32k (VeryLowLight-specific) and 84k (dark-augmented)
crops. Photometric and distance augmentation is applied identically to both
classes so that brightness and sharpness cannot serve as shortcuts.

Positives for the large sets were mined as the detector's top detection on
a gesture-labelled image, which confuses \emph{the image contains a hand}
with \emph{this detection is the hand}. An independent arbiter (the
landmark model's presence head, used offline) rejected 41.6\% of 30,000
mined crops, reproduced at 41.9\% on 31,006 further detections. Training on
the unfiltered set cut held-out AUC from 0.965 to 0.808.

\subsection{Motion gate and motion blobs}
\label{subsec:motion}
Every remaining false lock sat on a static object, while even a held-still
hand has micro-tremor. We apply a motion gate during acquisition only: over
a 5-frame, 320-px-wide luma ring, the fraction of pixels whose bias- and
gain-compensated value changes by more than $12/255$. Mean absolute
difference and temporal standard deviation both failed a deliberate test,
because they could not separate a held-still hand (score 0.036 to 0.083)
from the noisiest static clutter tile. The changed-pixel fraction separates
them with at least a 3.6$\times$ margin and is immune to a simulated
$\pm15\%$ auto-exposure ramp (0.0000 compensated against 1.0000
uncompensated). This allowed the acquisition size gate to drop from 0.15 to
0.10. Connected components of the same difference signal also propose
candidate regions (motion blobs), confirmed by the same verifier.

\subsection{Cursor logic}
The state machine alternates between REACQUIRE (full-frame detection plus
motion gate) and TRACK (a search window around the last position, which
gives the detector more effective pixels on a far hand). A plausibility
gate rejects jumps above 3 frame-widths per second. An EMA ($\tau=0.35$ s)
smooths positions before a cursor limited to 1 screen-width per second. The
dark variant adds 3-pass spatial consensus on acquisition and a TRACK hold:
a candidate more than 0.06 frame-widths from the committed position waits
for 3 agreeing candidates or 6 frames.

\section{Experimental Settings}

\subsection{Baselines}
MediaPipe Hands (0.10.14, tracking mode) is the incumbent and the method we
modify. A logistic regression on 8 interpretable crop statistics
(foreground fill, radial reach, compactness, edge density, intensity
moments; 9 parameters) is the simple learned baseline for the verifier. The
box-size gate is the ``before'' state for the verifier ablation.
Unmodified BlazePalm weights are the ``before'' state for the fine-tune.

\subsection{Data and splits}
Eight clips were recorded by the author: CloseLightMoving, CloseLightStill,
CloseDarkStill, FarLight, FarDark, VeryLowLight (651 frames, the target
condition), and two empty-room clips (NohandLight, NohandDark). Truth is
MediaPipe's static-mode output, lag-free, with gaps filled within $\pm5$
frames, giving 1,319 truth frames over the six hand clips (1,553 frames in
all). The first seven clips were recorded earlier and form the 7-clip
integrated evaluation used for ablations, which excludes VeryLowLight.
Precision, recall and F1 are per frame. Accepted frames are those where the
system reports a position.

Verifier training data come from HaGRIDv2 (11k to 130k crops; HaGRID-level
validation and test splits are held out) plus 100 crops from the author's
room. VeryLowLight is split into 50-frame blocks dealt to three roles by a
seeded shuffle, split seed 0: 250 training frames (blocks 0, 2, 3, 4, 5),
151 validation frames (blocks 7, 10, 11, 13) and 250 test frames (blocks 1,
6, 8, 9, 12). All model and threshold selection uses training and
validation data only, and the test blocks are scored once. Three limits
apply and are not resolved. (1) The validation blocks are not
representative of the test blocks, which is a property of the recording
rather than of the systems: brightness drifts across it, so blocks are not
exchangeable. Dark v3 scores F1 14.6 on the validation blocks and 43.5 on
the test blocks, a threefold difference, and we report both rather than the
flattering one. Threshold selection, which is what the validation blocks
are used for, is therefore weakly constrained (Section~\ref{sec:seeds}).
(2) The end-to-end F1 figures of Table~\ref{tab:main} are computed over the
whole VeryLowLight clip, which includes the blocks used for fine-tuning, so
that column flatters the fine-tuned configurations; on the test blocks
alone the same systems score MediaPipe 81.2, default 9.0, big-verifier 10.8
and dark v3 43.5. The other five hand clips were not used for fine-tuning,
though several operating parameters were chosen on them. (3) All dark data
is one person, room, lamp, camera and session.

\subsection{Parameters}
Table~\ref{tab:params} lists the key settings. Where we swept, we varied
one parameter at a time. Learning rates, batch sizes and epochs follow
common small-model defaults and were not swept.

\begin{table}[t]
\caption{Selected parameters and how they were chosen}
\label{tab:params}
\centering
\footnotesize
\begin{tabular}{@{}p{0.30\linewidth}p{0.13\linewidth}p{0.47\linewidth}@{}}
\toprule
Parameter & Value & Chosen by \\
\midrule
Acquisition size gate & 0.10 & swept 0.15, 0.10, 0.08, 0.06 \\
Motion gate floor & 0.01 & calibrated: 0.036 weakest still hand vs.\ 0.000 static room \\
Verifier thr.\ (default) & 0.30 & swept 0.20 to 0.50, clutter rejection at fixed recall \\
Verifier thr.\ (130k model) & 0.61 & swept, matched own-room clutter cut (91.3\%) \\
Detector positive weight & 5.0 & 20.0 collapsed the head \\
Detector matched anchors & 6 & about 96 (naive) failed \\
\bottomrule
\end{tabular}
\end{table}

The remaining settings come from measurement rather than sweeps: the
agreement rules use 3 passes within 0.05 frame widths to acquire and a TRACK
hold of 3 passes within 0.06 frame widths bounded to 6 frames, both read off
the measured frame-to-frame displacements; the plausibility limit is 3 frame
widths per second, the EMA $\tau=0.35$ s and the cursor cap 1 screen-width
per second are design choices from the hand-speed prior and a rate limit.

\subsection{Repeated runs, seeds and determinism}
All randomised components use a fixed, recorded seed list: split seed 0 for
the target-condition blocks, training seeds 0 to 9 for the verifier
variants and the detector fine-tune, and bootstrap seed 12{,}345 for the
resampling in Section~\ref{sec:results}. No seed is derived from clock
time. Section~\ref{sec:seeds} reports how many seeds each component
actually completed.

An earlier version of this harness gave F1 varying by about $\pm1.3$ points
between runs of the same fixed model, which we had called a noise floor. It
was not variance in the models. The plausibility gate (a speed limit in
frame widths per second) and the tracking loss timer are both time based,
and the replay harness fed them wall-clock \texttt{dt}, so replaying a clip
faster or slower than the machine happened to run changed which detections
were accepted. Timing now comes from the frame index and the clip's own
frame rate, and nothing in the decision path reads the clock. Repeating one
configuration three times then gives byte-identical output on all 16 clip
and pipeline rows, the only differing column being milliseconds per frame,
which measures the machine by design. The noise floor for evaluation is
therefore exactly zero, differences of any size are interpreted, and the
session caveat the previous comparison table carried is gone: all seven
systems are scored in one session against one truth. Mechanism changes are
checked separately by deterministic tests, a 13-sequence 6,300-step
equivalence test for a state-machine refactor and scripted flicker tests for
the consensus rules. Hardware: AMD Ryzen 5 7500F (6 cores), no GPU, Windows
10, Python 3.10.20, PyTorch 2.13.0+CPU, OpenCV 5.0.0.

\section{Results and Analysis}
\label{sec:results}

\subsection{Comparison with MediaPipe}
Table~\ref{tab:main} gives the main result. MediaPipe leads in aggregate
(F1 81.6\% against 45.4 to 52.5\%), mostly through recall (75.7\% against
32.7 to 40.0\%). At close range the gap narrows to 6 to 15 F1 points on the
first two clips, and no configuration produced a false lock on either
empty-room clip: every one of the seven systems scored 0.0\% on both. The
130k-crop verifier at 0.61 (52.5) is the best of our configurations on the
target clip, at more than twice the default's F1 there (10.4 to 21.5), and
it is also the best in aggregate, though Section~\ref{sec:stats} shows that
difference is not separable from noise.

\textbf{Contrast preprocessing does not help, and we test it rather than
assume it.} CLAHE in front of unmodified MediaPipe reaches F1 80.4 and
gamma 0.5 reaches 79.8, both below the 81.6 of raw MediaPipe, and both lose
more frames than they win: CLAHE is correct on 29 frames where raw
MediaPipe is wrong and wrong on 55 where raw is right (exact McNemar
$p_{\text{Holm}}=0.018$). The damage concentrates where the sensor is
already near its noise floor: on VeryLowLight both lose (65.1 raw, 61.8
CLAHE, 62.2 gamma), while on FarLight CLAHE gains slightly (75.0 to 76.7)
and gamma loses badly (66.7). CLAHE raises precision marginally (88.9
against 88.3) by accepting 42 fewer frames, which is what amplifying local
contrast on noisy footage does: more selective without being better. What
this rules out is that the target-condition failure is a contrast deficit
a preprocessing stage can repair, consistent with the attribution below,
where the loss happens before any appearance-based decision.

\begin{table}[t]
\caption{F1 (\%) per clip, then aggregate precision, recall and false
locks. All seven systems were scored in one session against identical
ground truth, so no column is from a different run. Higher F1 is better;
lower FL is better. MP = MediaPipe in tracking mode. +C = unmodified
MediaPipe behind CLAHE; gamma 0.5 is quoted in the text. Big = 130k-crop
verifier at 0.61. Dark v3 = fine-tuned detector plus acquisition
consensus; v4 adds the TRACK hold. FL = false-lock rate (\%) on
NohandLight/NohandDark; every configuration scored 0.0 on both.}
\label{tab:main}
\centering
\scriptsize
\setlength{\tabcolsep}{2.5pt}
\begin{tabular}{@{}lcccccc@{}}
\toprule
 & MP & +C & Default & Big & Dark v3 & Dark v4 \\
\midrule
CloseLightMoving & 100.0 & 100.0 & 86.4 & 85.8 & 93.6 & 90.2 \\
CloseLightStill & 100.0 & 100.0 & 92.5 & 92.5 & 54.6 & 51.6 \\
CloseDarkStill & 100.0 & 100.0 & 93.8 & 90.1 & 9.4 & 6.1 \\
FarLight & 75.0 & 76.7 & 32.3 & 34.5 & 50.2 & 34.6 \\
FarDark & 77.8 & 75.3 & 14.9 & 12.1 & 50.9 & 51.6 \\
VeryLowLight & 65.1 & 61.8 & 10.4 & 21.5 & 31.5 & 28.0 \\
\midrule
Aggregate F1 & 81.6 & 80.4 & 50.3 & 52.5 & 48.8 & 45.4 \\
Precision & 88.3 & 88.9 & 79.2 & 77.0 & 62.6 & 74.2 \\
Recall & 75.7 & 73.4 & 36.8 & 39.9 & 40.0 & 32.7 \\
Accepted (of 1,553) & 1,131 & 1,089 & 614 & 683 & 844 & 581 \\
FL Light/Dark & 0/0 & 0/0 & 0/0 & 0/0 & 0/0 & 0/0 \\
\bottomrule
\end{tabular}
\end{table}

\subsection{Where the target-condition loss occurs}
Of the 291 VeryLowLight frames where lag-free MediaPipe sees a hand (a
stricter count than the 486 truth frames used for F1), the detector fails
to produce a hand-containing box on 187 (64.3\%): 77 with no detection and
110 off the hand (Fig.~\ref{fig:funnel}). The size and motion gates then
remove 18 more. With the default verifier, 70 of the 86 crops that reach it
are rejected (81\%) and only 16 frames (5.5\%) survive. With the 130k
verifier, 13 are rejected (15\%) and 73 (25.1\%) survive. The detector
therefore accounts for 68\% of all lost frames under the default verifier
and 86\% (187 of 218) once the verifier is fixed.

\begin{figure}[t]
\centerline{\includegraphics[width=0.49\textwidth]{fig2_funnel.png}}
\caption{Where VeryLowLight hand frames are lost, for two verifier
configurations. Counts and share of the 291 frames.}
\label{fig:funnel}
\end{figure}

This is not a confidence problem. The share of frames with any detection is
59.8\% at every score threshold from 0.3 to 0.7, with median best score
0.76. The detector fires confidently, just often not on the hand. On
clean, well-lit HaGRID images its top detection is off the hand in about
42\% of cases, so the failure is not specific to our footage.

\subsection{Ablations}
Table~\ref{tab:ablation} gives operating points on the 7-clip integrated
evaluation.

\begin{table}[t]
\caption{Ablations on the 7-clip integrated evaluation (the five original
hand clips plus two empty-room clips): hand frames kept (more is better)
and clutter locks (fewer is better). One run per row.}
\label{tab:ablation}
\centering
\footnotesize
\begin{tabular}{@{}p{0.62\linewidth}rr@{}}
\toprule
Configuration & Hand frames & Clutter locks \\
\midrule
Box-size gate only (0.04) & 499 & 156 \\
Logistic baseline, 8 features (0.30) & 379 & 190 \\
CNN verifier at 0.20 & 490 & 143 \\
\textbf{CNN verifier at 0.30 (chosen)} & \textbf{486} & \textbf{32} \\
CNN verifier at 0.50 & 461 & 9 \\
+ motion gate, size gate 0.15 & 438 & 3 \\
+ motion gate, size gate 0.10 & 469 & 15 \\
\bottomrule
\end{tabular}
\end{table}

\textbf{Verifier.} At 0.30 the CNN keeps 97.4\% of true hand frames while
cutting clutter locks by 79\%. The logistic baseline is worse than no
filtering (it loses 120 hand frames and adds 34 clutter locks), and its
held-out AUC is 0.788 against 0.965 for the CNN. We do not have score-level
data for full ROC curves, so Table~\ref{tab:ablation} gives measured
operating points instead. The 130k verifier raised detector-recipe hand
recall on VeryLowLight crops from 4.9\% to 58.5\% at matched own-room
clutter rejection, but at lower portability (HaGRID clutter cut 91.5\% to
70.9\%, AUC 0.965 to 0.947).

\textbf{Motion gate.} Alone it removes 29 of 32 clutter locks but costs 48
hand frames, because it correctly refuses to bootstrap far-range
acquisitions from clutter. Relaxing the size gate to 0.10, safe only
because the motion gate now does the filtering, recovers 31 frames and cuts
clutter locks 53\% against the ungated verifier. Both empty-room clips give
zero acquisitions at every gate value down to 0.06.

\textbf{Motion blobs.} Across the 291 target-condition hand frames, blobs
propose the hand a little more often than the detector's box contains it
(36.8\% against 35.7\%), but on the frames where the detector fails outright
blobs propose a hand in 21.4\% of cases and the blob verifier confirms only
7.5\% of those, a ceiling of about 14 frames (4.8 points) before any
clutter cost. Without a blob
verifier they raise VeryLowLight clutter locks from 98 to 405 and lock onto
a synthetic moving non-hand object every time, so blobs are off in the
default. In the dark variant they stay on, because they enable
verifier-gated acquisition: turning them off raises the no-hand false-lock
rate from 1.0\% to 12.3\%.

\subsection{Detector fine-tune and consensus}
Fine-tuning gives the largest single gain, and it is a redistribution. With
the fine-tuned detector alone, VeryLowLight F1 rises 10.4 to 31.5 (28.0
with the TRACK hold, Table~\ref{tab:main}), FarDark 14.9 to 50.9 and
FarLight 32.3 to 50.2, but CloseDarkStill falls 93.8 to 9.4 and
CloseLightStill from 92.5 to 51.6. On the clean comparison, localisation on
VeryLowLight held-out blocks rises 25.0\% to 33.0\% (train-minus-holdout
gap 4.2 to 1.5 points) while CloseLightMoving falls 9.8 points. Since the
31.5 includes fine-tuning blocks, the held-out localisation gain is the
safer estimate of generalisation, and FarDark and FarLight are the better
evidence because those clips were not used for fine-tuning. In this
one-session evaluation the fine-tuned configurations produce no false locks
either: every system scored 0.0\% on both empty-room clips. An earlier
build of the dark variant measured 2.2\% false locks on NohandLight, and
that does not reproduce in the fixed-timing harness, which is reported here
because the number changed rather than because the new one is convenient.

The TRACK hold (Table~\ref{tab:main}, Dark v3 to v4) is measured with both
configurations in one process on the hand clips: off-hand cursor placements
fall 316 to 150, wrong-location excursions 79 to 30 (62\% fewer) and
persistent wrong locks 7 to 1, at a cost of 7.3 points of recall (40.0 to
32.7) and 3.4 of F1 (48.8 to 45.4). F1 counts frames with a reported
position, so holding the cursor lowers it by construction, and a cursor
that is briefly still is a different defect from one that jumps: F1 alone
undervalues this change and an excursion metric alone would hide the recall
cost. Acquisition consensus on its own is neutral end to end, and on
CloseDarkStill it fails structurally: false detections are stable and true
ones transient, so agreement over time discards the truth (F1 93.8 to 9.4).
That clip needs a crop-level discriminator, not more timing logic.

\subsection{Training runs across seeds}
\label{sec:seeds}

Ten seeds were planned for each of the four verifier variants and for the
detector fine-tune, and the campaign was stopped before it finished. We
report the counts as they stand rather than implying a uniform campaign:
ten seeds for the 11k variant, nine for the target-condition variant, three
for the 130k variant, none for the 84k dark-augmented variant, and five for
the detector fine-tune. The per-seed end-to-end pipeline evaluation, which
would turn these into per-clip and aggregate F1 means, was not run, so the
table reports what the seeds actually vary: the verifier's own held-out
AUC, the epoch its validation loss selected, and the threshold chosen for
it on the validation blocks.

\begin{table}[t]
\caption{Verifier variants retrained from scratch with the shared seed
list, mean $\pm$ standard deviation over the seeds that completed. HaGRID
AUC is on the HaGRID test crops (blank where the variant's pool contains
none); target AUC is on the VeryLowLight test blocks (blank where the pool
has no target-condition crops). Threshold is the per-variant operating
point selected on the 151 validation frames, and its spread is a property
of that validation set, not of the seed.}
\label{tab:seeds}
\centering
\footnotesize
\setlength{\tabcolsep}{3pt}
\begin{tabular}{@{}lcccccc@{}}
\toprule
Variant & Seeds & HaGRID AUC & Target AUC & Thr. & Epoch & Train s \\
\midrule
11k mined         & 10 & 0.975$\pm$0.005 & -- & 0.08$\pm$0.06 & 7.8$\pm$4.0 & 25 \\
130k arbitered    &  3 & 0.951$\pm$0.008 & 0.64$\pm$0.07 & 0.10$\pm$0.07 & 2.0$\pm$1.7 & 576 \\
Target-cond.      &  9 & 0.972$\pm$0.006 & 0.71$\pm$0.08 & 0.23$\pm$0.16 & 7.1$\pm$4.3 & 22 \\
84k dark-aug.     &  0 & \multicolumn{5}{c}{not run} \\
\bottomrule
\end{tabular}
\end{table}

Two findings come out of the seeds that did run, and they point in opposite
directions. The weights are stable: held-out AUC varies by 0.005 to 0.008
across seeds, a factor of three or four below the 0.02 AUC gap between the
130k variant and the two small ones, so that ranking is not seed noise. The
operating threshold is not stable: it ranges from 0.05 to 0.50 across the
seeds, standard deviation 0.06 to 0.16 within a variant, because 151
validation frames cannot pin it down. The detector fine-tune completed five
seeds: on the target-condition test blocks the share of frames whose
detection boxes contain the hand is 45.5\% (SD 2.0), every seed selected its
final epoch, mean training time was 222 s, and the collapse guard never
tripped, the worst seed leaving 5 anchors above 0.5 against a limit of 100.

\subsection{Domain dependence and statistical position}
\label{sec:stats}
The dark verifier keeps 61.3\% of held-out hand crops in the recorded
colour cast, 38.7\% under a warmer relighting, 52.4\% under a cooler one,
and only 4.0\% with the cast removed. It has partly learned this lamp's
colour. This is a measured generalisation failure, not a hypothetical one.

We can test, because the per-frame decision of all seven systems was kept
for every frame of every clip. Those decisions are paired binary outcomes on
the same frames, so the exact McNemar test on the discordant pairs answers
``does A decide better than B'' without discarding the pairing or assuming a
distribution. F1 is not a mean over independent frames, so an F1 difference
gets a 95\% interval from a block bootstrap over 50-frame blocks within each
clip (2{,}000 resamples, seed 12{,}345), which keeps the local correlation
of neighbouring frames. The ten comparisons form one family,
Holm-Bonferroni adjusted, $\alpha=0.05$.

\begin{table}[t]
\caption{Paired tests over the six hand clips (1,553 frames). A/B counts
frames where A decided correctly and B did not, and the reverse; $p$ is the
exact McNemar test on those discordant pairs, Holm-Bonferroni adjusted
across all ten comparisons. $\Delta$F1 carries a 95\% block-bootstrap
interval (50-frame blocks, 2{,}000 resamples, seed 12{,}345). A dagger
marks a frame-level difference that is significant while the F1 interval
still spans zero, so the two tests are answering different questions.}
\label{tab:stats}
\centering
\scriptsize
\setlength{\tabcolsep}{3pt}
\begin{tabular}{@{}llccccc@{}}
\toprule
A & B & A/B & $p$ (Holm) & $\Delta$F1 & 95\% CI & Sig. \\
\midrule
MP raw & default & 553/38 & 3e-117 & $+0.31$ & $+0.226,+0.403$ & yes \\
MP raw & big & 538/60 & 3e-96 & $+0.29$ & $+0.195,+0.384$ & yes \\
MP raw & dark v3 & 608/84 & 5e-98 & $+0.33$ & $+0.244,+0.414$ & yes \\
MP raw & dark v4 & 638/57 & 3e-124 & $+0.36$ & $+0.281,+0.445$ & yes \\
MP raw & MP CLAHE & 55/29 & 0.018 & $+0.01$ & $-0.002,+0.025$ & yes$^{\dagger}$ \\
MP raw & MP gamma & 66/30 & 2e-03 & $+0.02$ & $+0.002,+0.033$ & yes \\
MP CLAHE & MP gamma & 43/33 & 0.302 & $+0.01$ & $-0.009,+0.022$ & no \\
default & big & 54/91 & 0.011 & $-0.02$ & $-0.066,+0.014$ & yes$^{\dagger}$ \\
big & dark v3 & 262/216 & 0.079 & $+0.04$ & $-0.036,+0.103$ & no \\
dark v3 & dark v4 & 112/55 & 7e-05 & $+0.03$ & $+0.001,+0.067$ & yes \\
\bottomrule
\end{tabular}
\end{table}

The comparison the report rests on is supported without qualification: each
of our four configurations differs from MediaPipe at
$p_{\text{Holm}}<10^{-95}$, with an interval nowhere near zero. The rest of
the family is largely negative and is reported that way. Big against dark v3
is not significant ($p_{\text{Holm}}=0.079$), so its 3.7-point aggregate
advantage is not claimed; CLAHE against gamma is not significant either
($p_{\text{Holm}}=0.302$), so the two preprocessing variants are treated as
equally unhelpful rather than ranked. Two rows are significant at the frame
level while their F1 intervals span zero, and they are the ones to read
carefully: CLAHE loses more frames than it wins (55 against 29) but its F1
difference is inside the noise, so the claim is that it loses frames, not
that it lowers F1; and the retrained verifier decides correctly on 91 frames
where the default does not against 54 the other way
($p_{\text{Holm}}=0.011$) while its F1 difference of 0.02 is not separable
from zero, a real frame-level improvement that F1 alone would hide. The
41.6\% mined-label error rate keeps its binomial interval ($\pm0.6$ points,
ignoring clustering), and the arbiter ablation is re-measured on one
identical held-out pool of 4{,}071 crops: AUC 0.860 unfiltered against 0.959
arbitered, mean $P(\text{hand})$ on clutter 0.44 against 0.20.

\subsection{Computational cost}
BlazePalm has 1,763,358 parameters (7.1 MB). The verifier adds 60,929 (255
KB) and the fine-tune adds none. The pipeline runs at 19 to 22 ms per frame
on CPU with blobs off and 21 to 40 ms with blobs on, against 0.2 to 0.3 ms
for MediaPipe's cached detection step: about two orders of magnitude
slower, but still about 50 fps, enough for slow gestures. One anomalous run
measured the same configuration 2.2 to 4.6 times slower under host
contention, so timings are indicative. Verifier training took 22 to 576 s
per variant (mean over the seeds that completed, so the spread between
variants is the crop-set size), the detector fine-tune 222 s, and the
failed fine-tune 210 s. Building the system needed far more
than it ships: about 10 GB of HaGRIDv2 and 1.4 million generated crops for
a 255 KB verifier.

\section{Ethical Considerations}
\label{sec:ethics}

\textbf{Licensing and provenance.} HaGRID is released under CC BY-SA 4.0
\cite{hagrid}. We obtained HaGRIDv2 from two third-party mirrors
(\texttt{zhouxzh/portable-hagridv2-mediapipe-hand} and
\texttt{testdummyvt/hagRIDv2\_512px}) rather than the official release, so
completeness and licence-file preservation are not guaranteed. Whether a
model trained on CC BY-SA images is ``adapted material'' is unsettled, and
we state the licence without a legal conclusion. No HaGRID images appear
here.

\textbf{Privacy.} All eight clips show only the author's hand and room,
recorded with consent. HaGRID's subjects consented to that dataset's
release, which we rely on. The application processes frames in memory,
writes nothing to disk and makes no network calls.

\textbf{Bias and representation.} Dark training and test data come from one
person, room, lamp and camera (about 200 labelled frames), so no cross-user
or cross-environment generalisation is shown, and the colour-cast result
above shows the dependence. Label cleaning uses MediaPipe's landmark head
as arbiter, so any blind spot of that model (skin tone, hand size, unusual
posture) is inherited by our data and is not characterised.

\textbf{Safety and misuse.} Misuse potential is low: the output is local
cursor coordinates. The realistic risk is presenting this as an
accessibility aid. The system reports a position on only 37.9\% of frames
in its default configuration, has documented wrong-location excursions, and
was trained on one user, so it is a classroom prototype and not validated
for people who depend on it. Moving the cursor to a screen corner aborts
the program (pyautogui fail-safe) and is the only user-facing stop. The
motion gate assumes a stationary camera.

\textbf{Environmental cost.} Development used about 10 GB of downloads, 1.4
million generated crops and 16 PyInstaller builds. Energy was not measured.

\section{Conclusions and Future Work}

\textbf{Answer.} Partly, and only in specific conditions. MediaPipe wins in
aggregate (F1 81.6\% against 45.4 to 52.5\%), and paired tests confirm the
gap at $p_{\text{Holm}}<10^{-95}$. At close range the gap is 6 to 15 points
with identical false-lock behaviour, zero for every configuration. In the
target regime the dark variant raises VeryLowLight F1 from 10.4 to 31.5
(3.0$\times$; 1.5$\times$ over the 130k-verifier configuration at 21.5), but
still trails MediaPipe's 65.1 by 34 points, and the 31.5 is partly
in-sample because that column includes the blocks used for fine-tuning. The
gain is a redistribution, with a close-static collapse from 93.8 to 9.4. Two
interventions that looked plausible from the signal-processing side were
tested and failed: classical contrast enhancement in front of MediaPipe,
and motion-based proposals aimed at the detector's misses. The binding
constraint is the palm detector, which loses 64.3\% of target hand frames
before any added component acts. Our single narrow intervention there gave
the largest gain.

\textbf{Limitations.} One dark session, one person; the verifier variants
completed 3 to 10 seeds and the detector fine-tune 5, so the end-to-end F1
figures still come from single runs even though the verifier's own AUC has
a variance estimate; MediaPipe-derived ground truth, with no human labels,
so every number measures agreement with MediaPipe; validation blocks that
are not representative of the test blocks, and every threshold quoted is
one draw from a wide interval because of it; one sample per condition; a
cursor only, with no click vocabulary. The training runs were cheap (22 to
576 s per model on CPU) and were stopped by time, not by compute, which is
the weakest part of this evaluation.

\textbf{Future work.} (1) Finish the seed campaign, and run the per-seed
end-to-end pipeline evaluation so the end-to-end F1 figures carry a variance
estimate rather than only the verifier's AUC. (2) Select thresholds on more
validation frames, or select them per seed and report the spread, since the
present 151 frames leave the threshold uncertain by a factor of ten. (3)
Hand-label a subset of VeryLowLight to break the circular ground truth; the
150-frame list and the labelling tool exist but no labels were collected.
(4) Fine-tune the detector on broader artificial-light data, which is the
64.3\% loss and the only component with demonstrated headroom. (5) Record a
second dark session in a different room, and real moving clutter for the
blob path. (6) Add a crop-level discriminator for stable false positives,
which is what CloseDarkStill needs.

\section*{Statement}
This report and the project it describes are my own work. Code:
\url{https://github.com/Filkspline/AI339ProjectRepo}. Tools: Python,
PyTorch \cite{pytorch}, OpenCV \cite{opencv}, MediaPipe \cite{mediapipe_github},
pyautogui, PyInstaller, \LaTeX.

Generative AI: I used DeepSeek's agentic coding harness (DeepSeek V4),
run directly in the project folder, to set up the project's file structure
and to help implement and run the system under my own iterative, continuous
instruction. Across development I set the direction and the experiments to
run, reviewed every result it reported, and made every design decision,
threshold and scope decision in this report myself; the agent did not work
autonomously or choose what to build, measure, or report. I also used an
AI assistant (Claude) to help plan prompts for the coding agent and to help
draft and edit this report from my own project notes and the agent's
measured results. I reviewed and am responsible for all code, results,
claims and references in this report. External resources: BlazePalm
\cite{blazepalm_port} and its released weights \cite{mediapipe_github},
MediaPipe Hands \cite{mediapipe}, and HaGRIDv2 \cite{hagrid2} via the two
mirrors named in Section~\ref{sec:ethics}.

\begin{thebibliography}{00}
\bibitem{mediapipe} F.~Zhang, V.~Bazarevsky, A.~Vakunov, A.~Tkachenka,
G.~Sung, C.-L.~Chang, and M.~Grundmann, ``MediaPipe Hands: On-device
Real-time Hand Tracking,'' arXiv:2006.10214, 2020.
\bibitem{mediapipe_github} Google, ``MediaPipe,'' GitHub repository, 2020.
[Online]. Available: https://github.com/google/mediapipe
\bibitem{blazeface} V.~Bazarevsky, Y.~Kartynnik, A.~Vakunov, K.~Raveendran,
and M.~Grundmann, ``BlazeFace: Sub-millisecond Neural Face Detection on
Mobile GPUs,'' arXiv:1907.05047, 2019.
\bibitem{blazepalm_port} V.~Satija, ``BlazePalm: PyTorch and CoreML
implementation of MediaPipe Hands,'' GitHub repository. [Online].
Available: https://github.com/vidursatija/BlazePalm
\bibitem{hagrid} A.~Kapitanov, K.~Kvanchiani, A.~Nagaev, R.~Kraynov, and
A.~Makhliarchuk, ``HaGRID -- HAnd Gesture Recognition Image Dataset,'' in
\emph{Proc. IEEE/CVF WACV}, 2024; arXiv:2206.08219.
\bibitem{hagrid2} A.~Nuzhdin, A.~Nagaev, A.~Sautin, A.~Kapitanov, and
K.~Kvanchiani, ``HaGRIDv2: 1M Images for Static and Dynamic Hand Gesture
Recognition,'' in \emph{Proc. WSCG}, 2025.
\bibitem{clahe} K.~Zuiderveld, ``Contrast Limited Adaptive Histogram
Equalization,'' in \emph{Graphics Gems IV}. Academic Press, 1994, pp.
474--485.
\bibitem{pytorch} A.~Paszke et al., ``PyTorch: An Imperative Style,
High-performance Deep Learning Library,'' in \emph{Proc. NeurIPS}, 2019.
\bibitem{opencv} G.~Bradski, ``The OpenCV Library,'' \emph{Dr.\ Dobb's
Journal of Software Tools}, 2000.
\end{thebibliography}

\end{document}
